\section{Introduction}
\label{sec:intro}

Sign languages are natural languages with their own lexicon and grammar, used by millions of Deaf
and hard-of-hearing people. Automatic sign language processing covers tasks of increasing
difficulty: \emph{isolated} sign recognition (one sign per clip, closed vocabulary), continuous
recognition, and translation to and from spoken languages
\citep{bragg2019interdisciplinary,koller2020survey,yin2021including}. This report addresses
isolated recognition only. The distinction is not merely terminological: a classifier over a few
hundred glosses ignores the grammar of the language, its use of space and its non-manual markers,
and therefore cannot translate \citep{yin2021including}.

\emph{Pose-based} approaches first extract body, hand and face landmarks with an off-the-shelf
estimator such as MediaPipe Holistic
\citep{lugaresi2019mediapipe,bazarevsky2020blazepose,zhang2020mediapipehands,kartynnik2019facemesh},
then classify the resulting keypoint sequence. Compared with models operating on RGB video, they
are compact, largely insensitive to background and appearance, and allow inference without
transmitting video, which protects the privacy of users. The Google Isolated Sign Language
Recognition (ISLR) competition \citep{kaggle_islr} established this setting with
\numAslSamples{} landmark sequences of \numAslClasses{} American Sign Language (ASL) signs.

Three questions remain insufficiently addressed in this setting. They structure the report.
\begin{description}[leftmargin=1.2em,labelindent=0pt]
  \item[RQ1 Generalisation.] How accurately do pose-based models recognise signs produced by
        people absent from the training data, and how much does a random split overestimate this
        accuracy?
  \item[RQ2 Transfer.] Does pre-training on ASL, a comparatively well-resourced sign language,
        improve recognition in French Belgian Sign Language (LSFB), in particular when few
        labelled examples are available?
  \item[RQ3 Deployment.] Can the complete pipeline run in real time in a web browser on a laptop
        CPU while reproducing the predictions of the trained model?
\end{description}

\paragraph{Contributions.}
\begin{itemize}
  \item A single, unit-tested preprocessing function shared by training, evaluation and a
        JavaScript port, whose outputs match the reference implementation within \numJsParity{}
        (\cref{sec:method}).
  \item A signer-independent comparison of a statistical baseline, a bidirectional GRU and a
        convolution--Transformer on \numAslClasses{} ASL signs over three seeds, with paired
        significance tests, a measurement of split leakage and one-factor ablations
        (\cref{sec:results-asl,sec:ablations}).
  \item A controlled transfer study from ASL to LSFB, with learning curves from \numLcMin{}
        examples per sign to the full corpus and a linear probe of the pre-trained encoder
        (\cref{sec:transfer}).
  \item A fully client-side demo that classifies a sign in \numFpBrowser{}\,ms, and an analysis
        of INT8 quantisation, which first degrades accuracy because of feature outliers and,
        once corrected, remains slower on CPU (\cref{sec:efficiency}).
\end{itemize}

\Cref{sec:related} reviews related work. \Cref{sec:data,sec:method,sec:protocol} describe the
data, the method and the evaluation protocol. \Cref{sec:results} reports the results,
\cref{sec:discussion} discusses them together with their limitations, and
\cref{sec:conclusion} answers the three research questions.
